% Part: incompleteness
% Chapter: incompleteness-provability
% Section: first-incompleteness-thm

\documentclass[../../../include/open-logic-section]{subfiles}

\begin{document}

\olfileid{inc}{inp}{1in}

\olsection{第一不完全性定理}

我们现在可以描述 Gödel 最初对第一不完全性定理的证明。
设 $\Th{T}$ 为任意可计算公理化理论，其语言扩展了算术语言，且 $\Th{T}$ 包含 $\Th{Q}$ 的公理。
特别地，这意味着 $\Th{T}$ 能表示可计算函数和关系。

我们之前已经论证过，给定一种将公式和证明编码为数字的合理方式，
关系 $\Prf[\Th{T}](x,y)$ 是可计算的，其中
$\Prf[\Th{T}](x,y)$ 成立当且仅当 $x$ 是 $\Th{T}$ 中
Gödel数为 $y$ 的公式的某个推导的Gödel数。
事实上，对于 Gödel 所考虑的特定理论，Gödel 能够利用他论文中的
45 个函数和关系列表，证明这个关系是原始递归的。
对应于他所选择的特定的 $\Th{T}$，第 45 个关系 $x B y$ 正是 $\Prf[\Th{T}](x,y)$。
需要注意的是，在 Gödel 论文中用到“递归”一词之处，我们现在会使用“原始递归”这一术语。

由于 $\Prf[\Th{T}](x,y)$ 是可计算的，它在 $\Th{T}$ 中是可表示的。
我们将用 $\OPrf[\Th{T}](x,y)$ 表示代表它的公式。
令 $\OProv[\Th{T}](y)$ 为公式 $\lexists[x][\OPrf[\Th{T}](x,y)]$。
这描述了 Gödel 列表中的第 46 个关系 $\fn{Bew}(y)$。
如 Gödel 所指出的，这是唯一一个“不能被断言为递归”的关系。
他的意思大概是：从定义来看，无法直接看出它是可计算的；
而后来的发展实际上表明它是不可计算的。

设 $\Th{T}$ 为包含 $\Th{Q}$ 的可公理化理论。
则 $\Prf[\Th{T}](x, y)$ 是可判定的，
从而在 $\Th{Q}$ 中可由公式 $\OPrf[\Th{T}](x, y)$ 表示。
令 $\OProv[\Th{T}](y)$ 为我们上面描述的公式。
由不动点引理，存在公式 $!G_{\Th{T}}$，使得 $\Th{Q}$（从而 $\Th{T}$）推导
\begin{equation}
\ollabel{eqn:qpf}
!G_{\Th{T}} \liff \lnot \OProv[\Th{T}](\gn{!G_{\Th{T}}}).
\end{equation}
注意，$!G_{\Th{T}}$ 实质上断言了“$!G_{\Th{T}}$ 在 $\Th{T}$ 中不可推导”。

\begin{lem}\ollabel{lem:cons-G-unprov}
若 $\Th{T}$ 是一致、可公理化且扩展 $\Th{Q}$ 的理论，
则 $\Th{T} \Proves/ !G_{\Th{T}}$。
\end{lem}

\begin{proof}
假设 $\Th{T}$ 推导 $!G_{\Th{T}}$。
那么确实存在一个推导，因此对某个数 $m$，关系 $\Prf[\Th{T}](m,
\Gn{!G_{\Th{T}}})$ 成立。
于是 $\Th{Q}$ 推导语句
$\OPrf[\Th{T}](\num m, \gn{!G_{\Th{T}}})$。
所以 $\Th{Q}$ 推导
$\lexists[x][\OPrf[\Th{T}](x,\gn{!G_{\Th{T}}})]$，而根据定义，这即
$\OProv[\Th{T}](\gn{!G_{\Th{T}}})$。
由 \olref{eqn:qpf}，$\Th{Q}$ 推导 $\lnot !G_{\Th{T}}$，
并且由于 $\Th{T}$ 扩展 $\Th{Q}$，$\Th{T}$ 也推导 $\lnot !G_{\Th{T}}$。
我们证明了：如果 $\Th{T}$ 推导 $!G_{\Th{T}}$，那么它也推导 $\lnot !G_{\Th{T}}$，
从而它就会不一致。
\end{proof}

\begin{defn}
\ollabel{thm:oconsis-q}
一个理论 $\Th{T}$ 是\emph{$\omega$一致的}，如果满足以下条件：
如果 $\lexists[x][!A(x)]$ 是任意语句，并且 $\Th{T}$ 推导 $\lnot
!A(\num 0)$、$\lnot !A(\num 1)$、$\lnot !A(\num 2)$、……，
那么 $\Th{T}$ 不推导 $\lexists[x][!A(x)]$。
\end{defn}

注意，每个 $\omega$一致的理论也是一致的。
这只需由以下事实即可得出：如果 $\Th{T}$ 不一致，那么
对每个 $!A$ 都有 $\Th{T} \Proves !A$。
特别地，如果 $\Th{T}$ 不一致，它既推导对任意 $n$ 的 $\lnot !A(\num n)$，
也推导 $\lexists[x][!A(x)]$。
所以，如果 $\Th{T}$ 不一致，它就是 $\omega$不一致的。
由逆否律，如果 $\Th{T}$ 是 $\omega$一致的，则它必然是一致的。

\begin{lem}\ollabel{lem:omega-cons-G-unref}
若 $\Th{T}$ 是 $\omega$一致、可公理化且扩展 $\Th{Q}$ 的理论，
则 $\Th{T} \Proves/ \lnot !G_{\Th{T}}$。
\end{lem}

\begin{proof}
我们证明：如果 $\Th{T}$ 推导 $\lnot !G_{\Th{T}}$，那么它就是 $\omega$不一致的。
假设 $\Th{T}$ 推导 $\lnot !G_{\Th{T}}$。
如果 $\Th{T}$ 不一致，它已经是 $\omega$不一致的，得证。
否则，$\Th{T}$ 是一致的，于是由 \olref{lem:cons-G-unprov}，它不推导
$!G_{\Th{T}}$。
由于 $\Th{T}$ 中不存在 $!G_{\Th{T}}$ 的推导，$\Th{Q}$ 推导
\[
\lnot \OPrf[\Th{T}](\num 0, \gn{!G_{\Th{T}}}), \lnot \OPrf[\Th{T}](\num 1,
\gn{!G_{\Th{T}}}), \lnot \OPrf[\Th{T}](\num 2, \gn{!G_{\Th{T}}}), \dots
\]
并且 $\Th{T}$ 也推导它。
另一方面，由 \olref{eqn:qpf}，$\lnot
!G_{\Th{T}}$ 等价于
$\lexists[x][\OPrf[\Th{T}](x,\gn{!G_{\Th{T}}})]$。
因此 $\Th{T}$ 是 $\omega$不一致的。
\end{proof}

\begin{prob}
  每个 $\omega$一致的理论都是一致的。证明反之不成立，即存在一致但 $\omega$不一致的理论。
  通过证明 $\Th{Q} \cup
  \{\lnot !G_{\Th{Q}}\}$ 是一致但 $\omega$不一致的来完成此题。
\end{prob}

\begin{thm}
\ollabel{thm:first-incompleteness} 设 $\Th{T}$ 为任意
$\omega$一致、可公理化且扩展 $\Th{Q}$ 的理论。
则 $\Th{T}$ 不是完备的。
\end{thm}

\begin{proof}
  若 $\Th{T}$ 是 $\omega$一致的，则它是一致的，因此由 \olref{lem:cons-G-unprov} 可得 $\Th{T}
  \Proves/ !G_{\Th{T}}$。
  由 \olref{lem:omega-cons-G-unref} 可得 $\Th{T} \Proves/ \lnot !G_{\Th{T}}$。
  这意味着 $\Th{T}$ 是不完备的，因为它既不推导 $!G_{\Th{T}}$，也不推导 $\lnot !G_{\Th{T}}$。
\end{proof}

\end{document}